\section{Expected result and scope}

The expected advantage is not a stronger optimizer. It is a cleaner causal sequence:
\[
  \text{trajectory evidence}
  \longrightarrow
  \text{generated thought}
  \longrightarrow
  \text{later behavior}.
\]
OPSD may remain more efficient when the desired correction is only a local probability
shift and the intermediate thought is unimportant. Hinted SFT is more appropriate when
the diagnosis should become part of the model's ongoing reasoning and influence several
later decisions.

The main risk is that the model learns a generic sentence without using it. This is why
the hint-removal test matters. A second risk is training unsupported hints, which would
teach confident guesses. The dataset should therefore contain only hints justified by
the trajectory available before the hint.

The proposed contribution is a simple data and masking choice:

\begin{center}
  \textbf{Insert the desired realization as assistant thinking, train only those hint
  tokens, and mask everything else.}
\end{center}

At deployment, nothing special happens. The model generates its normal thinking; when it
encounters the learned situation, it should produce the corrective hint, retain it in the
autoregressive context, and continue accordingly.
